\section{交付了什么}

\textbf{军师。} 引擎枚举双方全部合法行动，对每一种组合跑一遍共享结算器，模型只负责把算好的事实讲成玩家能理解的取舍；双方同速时两种出手顺序各算一次，搜索不会偷偷利用随机数。局内主动提示也归军师而不是陪练，触发条件是代码里的五条：伙伴倒下、在两个选项间来回犹豫、真实枚举出来的明显策略错误、长停留，以及血量低于 35\% 而背包里还有药；「安静档」是玩家自己设的一个档位：这一档下教练不主动开口，只在玩家提问时回答——它不是关闭教练，而是把教练改成被动。它和「本局不再提醒」优先于上面任何一条触发条件。

\shot[0.94]{coach-evidence.png}

图中展开的那一块是这条提示的计算依据：双方当前血量与能量，以及它引用的两张知识卡。依据里的数字全部来自引擎枚举，模型只写最上面那句话。

\textbf{陪练。} 它的重点是陪伴。它记得住：状态从真实对战历史派生——连败几局、最近赢没赢、上一局打到第几回合、还剩几瓶药、本命是哪只、上次玩是几天前；这些字段决定它用哪一档语气，档位本身会作为证据写进模型输入，所以它说出的每句话都能回溯到一条本机记录。它在场：对局中间会自己插一句看法，不提问、不给指令。它有情绪，也会吐槽局面，比如「潮甲龟连着 2 个回合被草系按着打，我看得有点急」。

\shot[0.5]{companion-bubble.png}

气泡在战斗页左下角，独立于聊天面板：不打开聊天也能看见它，可以点掉，点掉之后本局静音。每局主动搭话有上限，最近七天的关闭次数会降低这个上限，四次之后直接归零。它不评价玩家水平、不给战术指令、不说教，这三条由 \code{checkCompanionRestraint} 在出句后扫描。

\subsection{老师}

老师不只是复盘。它盯着玩家看哪一个技能不动就讲哪一个：属性、能耗、算在当前目标身上多少伤害、打完剩几豆、什么条件下比手里别的选项更合适，但不催出招。每回合结束给这一回合的事实；整局结束给复盘，从减员和生命变化里挑三个回合逐个分析。小测按真实面板出变式题（「速度 38 加一次敏捷+3，对速度 41 的先手还是后手，还是无法确定」），选项里带「不确定」。培养建议是它唯一涉及成长的入口，给出加点收益与这次加点是否改变先手关系。

\shot[0.94]{teacher-review.png}

截图里是整局结束后的复盘：上面一行是这局的结论，展开区是它挑出来的关键回合，每个回合都带枚举算出的分数差。

真正花心思的是教学账本。它记三件事：讲过哪一课、哪一课被标回未掌握、这一课有几个被提示过和没被提示过的行动。讲过又没有新反证就不重复；只有军师在局内独立观察到三次同类失误（\code{RELEARN}：至少 3 次行动、合理率低于 50\%）才标回未掌握，并能说出「这一课教过，但之后 4 次独立行动里有 3 次不合理」。重新讲过就消费掉，再教要靠新出现的独立失误。

\shot[0.46]{growth-advice-table.png}

截图是培养建议：加 1 点力量之后攻击从 27 到 31，并摊开三项数值的对照。

这里有一个容易踩的坑，我是从文献里看到的。讲 assistance dilemma 的那篇论文报告，他们早期的预测器没把「学生用过提示」算进去，结果因为提示本身会把学生推向高效步骤，用过的学生反而被判成「不需要帮助」；他们的修正就是把两者分开。小芽里是同一件事：判断玩家是否掌握某个知识点时，只用他没有被提示过的时候做出的选择。

复盘一律使用回合开始时的公开状态，不使用对手实际打出的那一招——否则任何一次失败都能被解释成「你当时选错了」。改掉对手的实际行动，复盘结论不变，这条有测试。

这套账本是自建的计数与规则，不是心理测量模型：只判断教过没有、这次要不要再教，不估计掌握概率。

\subsection{游戏本体}

营地练级与配招、出征选关选人、局内逐回合出招，三段构成一局。顶部两个入口：训练按关卡挑战固定对手，对局里对手自由配队，也可以是同机分屏的两名真人，两边各有一条教练。

\shotpair{camp-train.png}{培养面板：经验、升级与三项加点}{camp-skills.png}{配招：6 选 4 与携带物}

出征分两步：先选关卡，再按选择顺序挑三只伙伴，阵容建议由教练给出。

\shot[0.94]{deploy-stages.png}

\shot[0.82]{deploy-roster.png}

\code{npm start}（Node.js 20+，无第三方依赖）之后浏览器打开本地端口即可。\textbf{不接模型也能完整游玩，规则建议、复盘、培养建议和小测都不依赖模型；接入模型之后，这些解释改由模型生成。}密钥只留在后端进程内存里，不落盘。

语音、联网 PVP、依赖真人被试的验收都没做完，逐条写在第五节。
